% Faithful formal-publication transcription: complete AppendicesA andB.
\section*{AppendixA: Proof of Lemma10}
We restate Lemma10.
\begin{authorclaim}{Lemma17}
Let $L\in\mathbb R^{m\times m}$ be an invertible matrix and $A\in\mathbb R^{m\times m}$ such that $\|L^{-1}-A\|\le\epsilon_0$. For any nonnegative integer $k$, define
\[R(A,L,k)=\sum_{i=0}^k(I-AL)^iA.\]
Then, $\|L^{-1}-R(A,L,k)\|\le\|L^{-1}\|\cdot\|L\|^{k+1}\cdot\epsilon_0^{k+1}$.
\end{authorclaim}
\begin{proof}
For any matrix $Z$, the matrices $I$ and $Z$ commute, and so by a straightforward induction,
\[I-\sum_{i=0}^k(I-Z)^iZ=(I-Z)^{k+1}.\]
In particular, for $Z=AL$,
\[I-R(A,L,k)\cdot L=(I-AL)^{k+1}.\]
Thus,
\begin{align*}
\|L^{-1}-R(A,L,k)\|&=\|(I-R(A,L,k)\cdot L)\cdot L^{-1}\|\\
&\le\|L^{-1}\|\cdot\|I-R(A,L,k)\cdot L\|\\
&\le\|L^{-1}\|\cdot\|I-AL\|^{k+1}\\
&=\|L^{-1}\|\cdot\|(L^{-1}-A)\cdot L\|^{k+1}\\
&\le\|L^{-1}\|\cdot\|L\|^{k+1}\cdot\epsilon_0^{k+1}.
\end{align*}
\end{proof}
\section*{AppendixB: The space complexity of some pseudorandom objects}
In this section we show how to achieve the space complexity declared in Theorem8 and Theorem9. For the INW generator we choose a specific implementation with a small space complexity. The constructions are well known, and the variant of INW we use was explored by [12]. We give it here for completeness.
\subsection*{B.1 Nisan's generator}
\paragraph{Proof sketch of Theorem8.}
We are given parameters $n,\Sigma,w,\epsilon$. We set $X=[A]$ for $A=O(nw\Sigma/\epsilon)$. We let $H$ be a 2-universal family of hash functions over $X$ where $|H|=A^2$ and $h(x)$, for $h\in H$ and $x\in X$, can be computed in space $O(\log\log|X|)$ (see [17,18]).
Nisan's generator interprets the seed as $y,h_1,\ldots,h_{\log n}$, where $y\in X$, and $h_1,\ldots,h_{\log n}\in H$. For $j\in[n]$, the $j$-th symbol in the output of the generator is $h_1^{b_1}(h_2^{b_2}\cdots h_{\log n}^{b_{\log n}}(y))$, where $(b_1,\ldots,b_{\log n})\in\{0,1\}^{\log n}$ is the binary representation of $j$, and $h^b$ is either $h$, if $b=1$, or the identity function, if $b=0$. Given $y,h_1,\ldots,h_{\log n},j=(b_1,\ldots,b_{\log n})$ we can compute the $j$-th output symbol in the following two alternative ways.
\begin{itemize}
\item We can successively compute $h_j^{b_j}\cdots h_{\log n}^{b_{\log n}}(y)$ for $j=\log n,\ldots,1$, each time keeping the current $X$-symbol. This takes
\[O\left(\log\frac{nw|\Sigma|}\epsilon+\log\log n+\log\log|X|\right)=O\left(\log\frac{nw|\Sigma|}\epsilon\right)\]
space.
\item Alternatively, we can do the above computation using composition of space bounded reductions, resulting in space complexity
\[O(\log n\cdot\log\log|X|)=O\left(\log n\cdot\log\log\frac{nw|\Sigma|}\epsilon\right).\]
\end{itemize}
\hfill$\square$
\subsection*{B.2 A high min-entropy extractor}
To apply INW, we need a space-efficient seeded extractor with a small entropy loss in the high min-entropy regime. Goldreich and Wigderson [11] gave such a construction utilizing a regular expander $G=(V,E)$ with a small normalized second eigenvalue. For our expander, we choose a Cayley graph over the commutative group $\mathbb Z_2^n$ with a generator set $S\subseteq\{0,1\}^n$ that is $\lambda$-biased. It is well known that $\operatorname{Cay}(\mathbb Z_2^n,S)$ has normalized second largest eigenvalue at most $\lambda$. For the $\lambda$-biased set we choose a construction from [2]. Altogether, this unfolds for the following.
\begin{itemize}
\item For the $\lambda$-biased set $S$, first pick $q$ to be the first power of two larger than $n/\lambda$. The set $S$ is of cardinality $q^2$. For every $\alpha,\beta\in\mathbb F_q$ there is an elements $s_{\alpha,\beta}\in\mathbb Z_2^n$ where $(s_{\alpha,\beta})_i=\langle\alpha^i,\beta\rangle$, such that multiplication is in $\mathbb F_q$ and the inner product is over $\mathbb Z_2$. [2] showed the set is $\lambda$-biased.
\item We let $G=(V,E)$ with $V=\mathbb Z_2^n$ and $(x,y)\in E$ iff $x+y\in S$. $G$ is a $\lambda$-expander.
\end{itemize}
The extractor $\mathrm{GW}:\{0,1\}^n\times[D]\to\{0,1\}^n$ is defined by letting $G(x,i)$ be the $i$-th neighbour of $x$ in the graph $G$.
\begin{authorclaim}{Claim18}
Let $0<\Delta<n$ and set $G$ and GW as above. Then, $\mathrm{GW}:\{0,1\}^n\times[D]\to\{0,1\}^n$ is a $(k=n-\Delta,\epsilon)$ extractor with seed length $d=O(\Delta+\log(n/\epsilon))$ and space complexity $O(\log n\cdot\log(\Delta+\log(n/\epsilon)))$.
\end{authorclaim}
\begin{proof}
For correctness, note that the expander mixing lemma shows that GW is an $(n-\Delta,\epsilon=O(2^{\Delta/2}\lambda))$ extractor.
\paragraph{Seed length.} The seed length of this extractor is $\log|S|=O(\log(n/\lambda))=O(\log(n2^\Delta/\epsilon))=O(\Delta+\log(n/\epsilon))$.
\paragraph{Space complexity.} The space complexity of computing $\mathrm{GW}(x,y)$ given $x$ and $y$, is the space needed to compute $s_y\in S$ from $y=(\alpha,\beta)\in\mathbb F_q^2$, plus the space needed to compute $x+s_y$. The dominating step in computing $s_y$ is computing $\alpha^i$ (for $i\le n$) which can be done in $O(\log n\log\log q)$ with space composition. Altogether, the space needed is
\[O(\log n\cdot\log\log(n/\lambda))=O\left(\log n\cdot\log\log\frac{n2^\Delta}\epsilon\right).\]
We note that Healy and Viola [12] gave an extremely efficient implementation of the above AGHP generator, yielding a better space complexity of $O(\log(n+\log q))$ to compute $\langle\alpha^i,\beta\rangle$. However, in our overall setting of parameters it will make negligible difference.
\end{proof}
We remark that by using expanders with better dependence between $D$ and $\lambda$, one can get $d=O(\Delta+\log(1/\epsilon))$, but here we care more about the space complexity, and $\log n$ factors are negligible for us.
\subsection*{B.3 The INW generator}
\paragraph{Proof sketch of Theorem9.}
We consider the INW generator [15] instantiated with extractors (as, e.g., in [21]). We are given parameters $n,\Sigma,w$, and $\epsilon=\epsilon_{\mathrm{INW}}$. We set parameters $\Delta=\log w+O(\log(n/\epsilon))$, and $d$ as the seed length for the extractor of Claim18 for length $n$, error $\epsilon_{\mathrm{Ext}}=\epsilon/n$ and $\Delta$. We let $s=\log|\Sigma|+\log n\cdot2d$ and we assume $s\le n$. We let $\ell_i=s-i\cdot\Delta$ for $0\le i\le n$.
Given a seed $x\in\{0,1\}^s$ we view the computation of $\mathrm{INW}(x)$ as a full binary tree of depth $\log n$. Nodes in level $i$ of the tree are labeled by strings of length $\ell_i$. The root (at level 0) is labeled by $x$ (of length $\ell_0=s$). Given any internal node in level $i\in\{0,\ldots,\log n\}$ labeled by some string $z\in\{0,1\}^{\ell_i}$, we write $z=z_1\circ z_2$ with $z_i\in\{0,1\}^{\ell_{i+1}}$ and $z_2\in\{0,1\}^d$. The left child of $z$ is labeled with $z_1$, and the right child of $z$ is labeled with $\mathrm{Ext}_i(z_1,z_2)$, where $\mathrm{Ext}_i$ is given by Claim18 for $\Delta$, length $\ell_{i+1}$ and error $\epsilon_{\mathrm{Ext}}$ (notice that since $\ell_i<n$, $d$ bits suffice for the seed). $\mathrm{INW}(x)$ is the concatenation of the leaf's labels, from left to right, truncating outputs to $\log|\Sigma|$ bits.
Given an index $j\in[n]$, computing $\mathrm{INW}(x)_j\in\Sigma$ can be done by walking down the computation tree, and each time either truncating a string or invoking an extractor. By composition of space bounded reductions the space complexity of the construction is $\log n$ times the space complexity of the worst extractor used. That is, $\log n\cdot\log\ell_0\cdot\log(\Delta+\log(\ell_0/\epsilon_{\mathrm{Ext}}))$. Plugging-in $\Delta$ and $\epsilon_{\mathrm{Ext}}$, the space complexity is bounded by
\begin{align*}
O\left(\log n\cdot\log\ell_0\cdot\log\log\frac{nw}\epsilon\right)
&=O\left(\log n\cdot\log\left(\log|\Sigma|+\log n\log\frac{nw}\epsilon\right)\cdot\log\log\frac{nw}\epsilon\right)\\
&=O\left(\log n\cdot\left[\log\log\frac{nw|\Sigma|}\epsilon\right]^2\right).
\end{align*}
\hfill$\square$
